
%%%%%%%%%%%%%%%%%%%%%%%%%%%%%%%%%%%%
% Recap/Agenda
%%%%%%%%%%%%%%%%%%%%%%%%%%%%%%%%%%%%
% TODO better formatting
\begin{frame}
    \frametitle{Module 3: Foundations for Inference}
    \begin{itemize}
        \item \hl{Previously: }Point estimates and sampling variability, continued (Chapter 5.1)
        \item \hl{This time: }Confidence intervals (Chapter 5.2)
        \item \hl{Reading: }Chapter 5.3
    \end{itemize}

\end{frame}
%%%%%%%%%%%%%%%%%%%%%%%%%%%%%%%%%%%%
% Learning objectives:
%%%%%%%%%%%%%%%%%%%%%%%%%%%%%%%%%%%%
\begin{frame}
    \frametitle{Learning Objectives}
    \begin{itemize}
        \item \textbf{M3, LO2: Visualize and Interpret Sampling Distributions:} Draw and interpret sampling distributions for a point estimate (e.g., population proportion) across different sample sizes, explaining how the distribution changes as the sample size increases.
        \item \textbf{M3, LO3: Inference for a Single Proportion:} Design, execute, and interpret confidence intervals and hypothesis tests for a single proportion.
    \end{itemize}
\end{frame}
